\section{Every control on screen, by its exact name}
\hypertarget{ui-reference}{}

This chapter lists each visible button, selector, and panel using the words
printed on screen. Bold text is always an exact UI label. Use it as a map when
a label is unfamiliar; hovering over most controls in the app also shows a
one-line tooltip.

\newcommand{\UiTable}[1]{%
\begin{center}
\begin{tabularx}{\textwidth}{>{\bfseries\raggedright\arraybackslash}p{0.27\textwidth}X}
\toprule
\normalfont\bfseries On screen & \bfseries What it does\\
\midrule
#1
\bottomrule
\end{tabularx}
\end{center}}

\subsection{Navigation menu}
\hypertarget{ui-navigation}{}

The three-line button in the top corner (\emph{Open site navigation}) opens
the menu. Every page listed here is part of the same browser session: nothing
is sent to a server for simulation.

\UiTable{
Circuit builder & The main workbench: palette, canvas, run controls,
examples, and protocol compiler.\\
Wiki / docs & Short in-app explanations of circuit construction, reset
buttons, the learning path, and custom gates.\\
QPU Documentation & Opens these PDF guides inside the app.\\
Particle visualization & Per-qubit spheres, Bloch coordinates, the execution
timeline, final state vector, and execution log.\\
Circuit correction lab & Tests a protocol against a truth table and helps
repair it.\\
Upload files / Download files & Both open the file page for moving
\code{.qpucir} and \code{.qpuio} files.\\
More & Quick links, including \textbf{Reset site completely}.\\
Reset site & Returns the whole builder to its first-visit state. See
\hyperlink{reset-buttons}{Reset buttons compared}.\\
Docs on hover & The switch at the end of this row (also at the top of the
Circuit builder) opens a documentation drawer at the side of the screen.\\
}

\subsection{Docs on hover drawer}
\hypertarget{ui-docs-drawer}{}

With \textbf{Docs on hover} switched on, rest the pointer on an item (or reach
it with the keyboard) and its documentation appears in the drawer. Moving
quickly across other items on the way to the drawer does not replace it.

\UiTable{
Gate blocks and canvas gates & Summary, rule, truth table, what happens to
the target t, and the protocol syntax. On the canvas the drawer also says
which role that wire plays in that step.\\
Custom gates & The same, built from the saved process: a simulated truth
table, which wire each input and output uses, and the \code{RUNCHILD} line.\\
Bundled protocol buttons & The process's truth table, inputs, outputs, and
child processes.\\
Builder controls & What the control does and where it is documented.\\
Hovered & The tab that follows the pointer.\\
Keep in a tab & Copies the current entry into its own tab so it stays put.\\
Open \dots{} in a drawer tab & Opens a child process in a new drawer tab,
not a new browser tab. Children of children open the same way.\\
}

\subsection{Gate palette: ``Pick up a block''}
\hypertarget{ui-palette}{}

Click a block to select it, or drag it onto a wire of the circuit canvas. The
label printed on each block is shorter than the protocol instruction name:

\begin{center}
\begin{tabular}{llll}
\toprule
Block & Instruction & Block & Instruction\\
\midrule
X, Y, Z, H, S, T & same name & CX & \code{CNOT}\\
P & \code{PHASE} & CCX & \code{CCNOT} (Toffoli)\\
CZ, CY & same name & SW & \code{SWAP}\\
M & \code{MEASURE} & $\neg$ & \code{NOT}\\
AND, NAND, OR, XOR & same name & & \\
\bottomrule
\end{tabular}
\end{center}

\textbf{Preconfigured} holds the built-in gates. \textbf{Custom gates} holds
gates you registered from a process; it is empty until you register one.

\subsection{Gate engine: ``Register custom gates from QPU processes''}
\hypertarget{ui-custom-gates}{}

\UiTable{
Gate id (palette label) & The name of the new block. Must start with a letter
and cannot reuse a built-in gate name. Leave empty to use the process name.\\
Source process & Which process becomes the gate: \emph{Current protocol
editor} or any process in the catalog.\\
Color (optional --- random if empty) & Block colour, as a CSS colour or
gradient.\\
Register custom gate & Compiles the process and adds the block to
\textbf{Custom gates}.\\
Remove & Deletes that custom gate from this browser session.\\
}
See \emph{Custom gates} in Examples and Troubleshooting for how inputs and
outputs map onto wires.

\subsection{Circuit canvas: ``Standard circuit diagram''}
\hypertarget{ui-canvas}{}

Each row is one wire, labelled \textbf{q0}, \textbf{q1}, \ldots{} with its
start state. Time runs left to right: every gate occupies its own column.
\begin{itemize}
  \item \textbf{Placing a gate:} drop a block on a row, or select a block and
  click the row. The row you drop on becomes the \emph{target}. Controls are
  chosen automatically from the lowest other rows; use the workbench to choose
  them yourself.
  \item \textbf{Removing a gate:} click the gate symbol on its target wire.
  \item \textbf{Line styles:} a double line is a classical bit (a
  \code{0p}/\code{1p} start or a measured wire); a single line may hold a
  superposition. A red outline marks the step that just ran.
\end{itemize}

\subsection{Interactive workbench: ``Add particles, gates, and measurements''}
\hypertarget{ui-workbench}{}

The workbench places a gate with exactly the wires you choose.

\UiTable{
Gate & Which gate \textbf{Add gate to target} will place. Same list as the
palette.\\
Target particle & The wire the gate \emph{changes}. For a measurement, the
wire \textbf{Measure target} reads.\\
Control A & The first wire the gate \emph{reads}. Used by CNOT, CZ, CY, and as
the first input of CCNOT, AND, NAND, OR, and XOR.\\
Control B & The second wire the gate reads. Used by two-input gates (CCNOT,
AND, NAND, OR, XOR) and custom gates with two or more inputs. SWAP also uses
it as the wire to exchange with the target.\\
Phase angle & Appears for PHASE only. Sets $\theta$ in degrees.\\
Add gate to target & Appends the selected gate as the next column.\\
Add particle / Remove particle & Adds or removes the last wire (up to six on
the canvas). After a compiled protocol, adds or removes a \code{PARAMS}
input instead; compile again afterwards.\\
Measure target & Measures only the selected target wire now, without adding
an \textbf{M} block to the circuit.\\
\emph{name} start & One selector per input wire: \code{0p} ($\ket0$),
\code{1p} ($\ket1$), or \code{sp} ($\ket+$). Changing it resets the run.\\
}

Gates that do not use a selector ignore it. Single-qubit gates use only
\textbf{Target particle}. A control can never equal the target, so the app
disables that choice and picks a free wire if needed.

\begin{center}
\WorkbenchControlMap{CCNOT}
\end{center}
\begin{beginnerbox}[Reading the picture]
The selectors on the left produce the diagram on the right. \textbf{Control A}
and \textbf{Control B} become the two filled dots: they are read and never
changed. \textbf{Target particle} becomes the $\oplus$: it is the only wire
that can change, and it changes by $t' = t\oplus(A\land B)$. The same mapping
applies to AND. For NAND, OR, and XOR only the rule written on the target
changes.
\end{beginnerbox}

\subsection{Run controls}
\hypertarget{ui-run-controls}{}

\UiTable{
Play Sequence / Pause sequence & Runs one gate at a time at the chosen speed,
so you can watch each step. At the end it replays from the start.\\
Run all & Jumps straight to the final state. Useful once you trust the
circuit.\\
Step gate & Runs exactly one more gate. Disabled at the end of the
circuit.\\
Reset state & Rewinds the run to the start states and clears measurements.
Keeps every gate.\\
Measure all & Measures every unmeasured wire now.\\
Clear circuit & Deletes every gate. Keeps the wires and their start states.\\
Reset site & Returns the entire builder to its default state.\\
Speed & Playback speed for \textbf{Play Sequence}.\\
}
The difference between the three reset-like buttons, and between them and the
compiler's internal RESET operation, is explained in
\hyperlink{reset-buttons}{Reset buttons compared}.

\subsection{Examples: ``Load a starter circuit''}
\hypertarget{ui-examples}{}

Each card loads a small circuit onto the canvas, replacing the current gates.
Cards are labelled with their step in the learning path, so they can be worked
through in order.

\subsection{QPU protocol compiler: ``Compile text into a circuit''}
\hypertarget{ui-compiler}{}

\UiTable{
SingleBitFullAdder, TwoBitFullAdder, FourBitFullAdder, PhaseDemo & Load and
compile a bundled protocol in one click.\\
Protocol editor & The text box. Editing it does not change the canvas until
you compile.\\
Compile protocol & Turns the editor text into canvas gates and resets the run.
The sentence beside the buttons summarizes the result or the error.\\
Download as .qpucir / Download as -qpucir.txt & Saves the editor text. The
\code{.txt} form is for devices that cannot open custom extensions.\\
Protocol checks & Live error and warning list for the editor text, with line
numbers and suggested fixes.\\
Open Correction Lab for guided help & Appears when checks find a problem.\\
Supported operations and token map & Expands to show every accepted
instruction and which wire each token name uses.\\
}

\begin{cautionbox}[Canvas edits rewrite the editor]
Placing, removing, or clearing gates on the canvas regenerates the protocol
editor text from the canvas. Download a hand-written protocol before editing
its compiled circuit visually.
\end{cautionbox}

\subsection{Particle visualization}
\hypertarget{ui-particles}{}

\UiTable{
Qubit states and execution progress & One card per displayed wire: a sphere
while unmeasured, the measured bit afterwards, plus its ket, $r,\theta,\phi$,
and Bloch $x,y,z$. $\Delta$ values show the change caused by the last step.\\
Circuit execution timeline & One tile per gate, with its \code{-I} and
\code{-O} wires.\\
Final state vector & Every nonzero amplitude with its basis label.\\
Outcome truth table / probabilities & Expands to list every basis outcome
with its amplitude and probability.\\
Execution log & A sentence for every step, compile, and measurement.\\
}

\subsection{File upload and download}
\hypertarget{ui-files}{}

\UiTable{
Upload files & Choose a \code{.qpucir} file, optionally together with its
matching \code{.qpuio} truth table, in one dialog. The protocol compiles
immediately.\\
Download \emph{process} & Saves a bundled protocol exactly as shipped.\\
Download current editor protocol & Saves whatever is in the protocol editor,
even if it has not compiled.\\
}

\subsection{Circuit correction lab}
\hypertarget{ui-correction-lab}{}

\UiTable{
Upload protocol / Upload truth table & Load a \code{.qpucir} and, optionally,
a \code{.qpuio} table.\\
Catalog process & Pick any compiled, uploaded, or bundled process.\\
Infer truth table & Build an empty table whose columns match \code{PARAMS}
and \code{RETURNVALS}.\\
Full-adder table & Load the canonical one-bit full-adder expectations.\\
Probe outputs & Fill the output columns with what the current circuit
actually produces. Refused for protected bundled tables.\\
Input columns / Output columns & Resize the table.\\
Test circuit & Compare the circuit with every table row; failing rows are
highlighted.\\
Test \& correct automatically & Test, then propose gate-level fixes for
failing rows.\\
Download JSON / .qpuio / -qpuio.txt / .qpucir / -qpucir.txt & Save the table
or protocol.\\
AI model settings & Choose the browser model or an Ollama server for the
chat.\\
Use AI for unrecognized messages & Lets the model interpret requests the
built-in command patterns do not recognize.\\
Send correction & Sends the chat message, for example
``test the circuit''.\\
}
